Dear Prof. Rajendran and Members of the Selection Committee,

I am writing to apply for the \textit{PhD Studentship in Algorithms/Hardware Co-Design for Bayesian and Ising Models} (Reference: R143025) at Northeastern University London in collaboration with IMEC.

I hold a Bachelor of Technology in Computer Science with a specialisation in Artificial Intelligence from Amrita Vishwa Vidyapeetham. I currently work as a Research Assistant at the SustainAI Lab at the Indian Institute of Technology Guwahati, investigating spiking state space models and their practical applications. My research in hardware aware machine learning and physical neuromorphic computing has led me to explore how the properties of a computational system can be incorporated into algorithm design to improve computational efficiency. This perspective is what draws me to this project.

My work on \textit{SlimEdge} has given me direct experience in designing algorithms around hardware constraints. I developed a framework for optimising distributed neural network inference across heterogeneous edge devices, combining model pruning, performance estimation, and multi-objective optimisation to identify suitable deployments under latency, accuracy, memory, and device availability constraints. Evaluating thousands of deployment configurations showed me how hardware characteristics shape trade offs between accuracy, latency, and resource consumption. This experience is directly relevant to the project’s objective of jointly optimising probabilistic algorithms and hardware architectures. I would bring experience in formulating constrained optimisation problems, evaluating competing objectives, and designing systematic experiments to assess computational trade offs.

My research on neuromorphic computing has approached the same challenge from the perspective of physical computation. I developed a reservoir model based on nonlinear oscillator dynamics and studied how its operating conditions affected temporal prediction performance. This research resulted in a paper accepted for publication in \textit{Applied Physics Letters}. Rather than treating the physical substrate simply as a platform for computation, my research examined how its intrinsic dynamics can contribute to computational performance. I see a compelling connection with the proposed research on stochastic MRAM devices. Although the underlying physical mechanisms differ, both research directions investigate how the properties of a physical system can be exploited for computation. My experience studying the relationship between physical behaviour and computational performance provides a relevant foundation for this work.

What particularly attracts me to this studentship is the opportunity to bring these two perspectives together. My experience in hardware aware optimisation provides a foundation for investigating the algorithmic trade offs involved in probabilistic computing, while my work on physical neuromorphic computing has taught me to consider how physical dynamics can shape computational capabilities. I am keen to extend this experience to Bayesian and Ising models and investigate how MRAM device statistics can be exploited through joint algorithm and hardware optimisation. The project’s emphasis on connecting device level behaviour with algorithm design and evaluating system performance makes it a strong fit for my research interests.

I would welcome the opportunity to contribute to the collaborative research at Northeastern University London, University of Kent, and IMEC. The project would also allow me to deepen my expertise in probabilistic modelling and emerging memory technologies. My experience in hardware constrained optimisation and physical computing provides a strong foundation for this work.

Thank you for considering my application.